% Overview by GPT-6 Astra (OpenAI, ChatGPT Pro mode).
\begingroup\selectlanguage{english}
\noindent\textbf{From replacements to a calculus of coherent squares.}
The pages gathered here form a single mathematical movement. Grothendieck first constructs categorical cylinders and uses them to establish the gluing rules for homotopy pushouts. He then asks which part of this geometry can be recovered from a category $\mathcal M$ and a class $W$ of weak equivalences alone. The answer takes the form of a homotopy completion of a cospan, characterized by an adjunction. Finally he returns to pushouts and investigates how far weak equivalences can be recovered from the pieces of a decomposition. The first two movements build the calculus; the appendix tests a considerably stronger conservativity property.

These pages are manuscript pages 228--275, printed pages 145--178 of the edition \citep{Grothendieck1990XIIc}: the second part of section~VI of Chapter~XII, that is, the end of its section~6 followed by sections~7 and~8. The audit covers that entire range, including the fragment of a footnote continued at the beginning. Page labels in the margins reproduce the two numbering systems of the transcription; citations in the commentary use manuscript pages unless another system is specified.

\medskip\noindent\textbf{Why the diagrams must precede localization.}
The motivation is older than the technical apparatus of these pages. In section 69 of \emph{Pursuing Stacks}, Grothendieck traces the need for derivators to the trace formalism of SGA~5: a single derived category did not retain the diagrammatic information needed for the desired compatibilities \citep[\S69, pp.~1--2]{Grothendieck1983}. His letter to Thomason of 2 April 1991 returns to this history and to the system of localized diagram categories \citep[\S2, pp.~5--6]{Grothendieck1991Letter}. These are historical statements checked in those texts; we do not infer that every later definition was already present there in its final form.

Here this motivation becomes a concrete distinction. With the manuscript's contravariant convention, the relevant category is
\[
 \mathbb D(I)=\operatorname{Fun}(I^{\mathrm{op}},\mathcal M)
                  [W(I)^{-1}],
\]
where $W(I)$ consists of the objectwise weak equivalences. It is not obtained by first passing to $\mathcal M[W^{-1}]$ and then forming its ordinary diagram category. A cospan remembers two maps to a common base together; its homotopy completion must preserve that relationship and its naturality in all three vertices. This explains both the attention given to the whole completed square and the two localization lemmas whose proofs the manuscript postpones. The modern introduction in \citet[pp.~313--315]{Groth2013} makes the same structural distinction.

\medskip\noindent\textbf{The constructive part and its later justification.}
The fence interval $\Phi_n$, with $n\geq2$, lets both endpoints enter by a Dwyer immersion. The associated mapping cylinder gives a functorial factorization into a $W$-cofibration followed by a weak equivalence. The commentary proves the endpoint condition and follows the construction on the map from two discrete objects to a point. Taking its categorical cofiber produces an eight-edge circle. Later, the Grothendieck integral of the same span gives a four-edge circle. These are two explicit models of the homotopy pushout; the ordinary categorical pushout is only a point.

The appeal to Thomason is valid in the needed cellular form. It must not be strengthened to the assertion that every Thomason cofibration is a Dwyer map: the retract assertion used to obtain that conclusion is false, and \citet{Cisinski1999Dwyer} gives the counterexample and repair. The recent comparison of Dwyer pushouts with $(\infty,1)$-categorical pushouts is stronger in a different direction, identifying what the nerve preserves rather than adding the appendix's conservative axiom \citep[Theorems~1.5--1.6, arXiv v3]{HackneyOzornovaRiehlRovelli2024}.

In section 7, a $W$-fibration means a map whose pullback functor preserves the designated weak equivalences. The factorization assumption says that every arrow is a weak equivalence followed by such a fibration. It is essential to retain this intrinsic meaning: the manuscript has not assumed a Quillen model structure with arbitrary lifting classes. Under its stated terminal-object and pullback hypotheses, these data form a left derivable category. The commentary verifies the nontrivial axiom, stability of trivial fibrations under arbitrary base change, by a two-step factorization and two-out-of-three argument.

This observation gives a complete modern repair of the deferred Lemmas 1 and 2. We construct the appropriate fibration structures on the good cospans and good squares, check the factorization and pullback conditions, and verify approximation for their inclusions in the ambient diagram categories. Cisinski's Proposition~1.8 also supplies the strengthened replacement assertion for every full intermediate subcategory of a slice containing its fibrant objects; his Theorem~3.12 establishes the two diagram-localization equivalences \citep{Cisinski2010}. No functorial factorization or global assumption that all objects are fibrant is required. A separate, shorter deformation argument explains what functorial factorization would give. The audit still labels the source proofs as deferred or sketched: the later completion does not change what is written on these pages.

With those equivalences available, restriction from squares to cospans has a fully faithful right adjoint. On a cospan $A\to S\leftarrow B$ with one good leg, its missing corner is represented by $A\times_S B$. The groupoid example $*\to BG\leftarrow *$ gives the discrete groupoid $G$ at that corner, whereas the strict pullback without replacement is a point. The adjunction retains the two projections and their compatibility, as well as the corner's homotopy type. Proposition 19 then gives an intrinsic test for cartesianness, and we prove its independence of choices, invariance under objectwise weak equivalences, and pasting directly from the factorization axiom. Strong saturation enters precisely when an isomorphism after localization must be recognized as an original member of $W$.

\medskip\noindent\textbf{What the appendix adds, and where its generalization fails.}
The ordinary gluing implication says that equivalences on the common part and the two pieces induce an equivalence on their homotopy pushouts. The appendix asks for converses. Its full W(7 bis) hypothesis is substantially stronger: a homotopy pushout must also reflect weak equivalences in an attached leg. In a coherent homotopy-pushout theory this is equivalent to detection by the terminal homotopy cofiber
\[
 C_e(u)=e\amalg^{\mathrm h}_{\operatorname{dom}(u)}
                         \operatorname{cod}(u).
\]
The forward implication $u\in W\Rightarrow C_e(u)\simeq e$ is formal. The reverse implication is the additional detection property. We prove the equivalence with conservativity by pasting cofiber squares.

This distinction exposes an actual error in the general theorem on manuscript page 274. Retract closure of aspherical objects alone does not imply that equivalences at the initial and final vertices of a map of homotopy-pushout squares force equivalences at the two middle vertices. Our counterexample uses Hatcher's finite acyclic noncontractible complex $A$. Its suspension is contractible; the map from the square $(A,A,CA,CA)$ to $(A,*,CA,CA/A)$ has three weak-equivalence components and one component $A\to *$ that is not a weak equivalence. Both squares are homotopy pushouts. We give the presentation of $A$, its cellular homology calculation, a nontrivial permutation representation of its fundamental group, and the suspension argument \citep[Example~2.38]{Hatcher2002}. The related equivalence asserted in Corollary 1 therefore fails as printed. Adding cofiber detection repairs both statements, and the final Corollary 2 already has the required conservative hypothesis.

There is also a source slip in Cons~2: the edition prints \emph{cart\'esien} where its argument requires \emph{cocart\'esien}. A finite chain-complex example shows that the literal cartesian assertion fails even in a stable theory. The proof of the intended criterion shows that its cocartesian version follows from Cons~1 under the coherent formulation. In stable derivators the entire corrected cube calculus has a familiar expression through the Mayer--Vietoris triangle \citep[Theorems~6.1 and 7.1, arXiv v2]{GrothPontoShulman2014}.

The appendix's earlier circularity is a different issue. Chapter XVI identifies the dependence on the Dwyer-cofibration gluing lemma but explicitly does not give its proof there \citep[\S3, ms.~pp.~18--23]{GrothendieckDerivateursXVI}. Brown's published gluing lemma, dualized, supplies the missing result \citep[Part~I, \S4, Lemma~2]{Brown1973}. Thus the initial conditional proposition survives, while the later hypothesis-free generalization does not. These are distinct diagnoses, with distinct repairs.

\medskip\noindent\textbf{Editorial and verification limits.}
The text preserves slips printed in the edition. Its one mathematical transcription correction restores the interior natural transformation $\alpha$ on manuscript page 235, which the transcription had drawn as an arrow between the wrong objects. The separate registers record the source's proof status, the present repair, and the evidence for each reading. Bibliography notes identify exactly which statements, proofs, or metadata were checked; in particular, the original proof of Thomason's 1979 theorem was not inspected, although its statement and a later proof were checked.

No assertion is labelled currently open merely because its proof is absent from these pages. The constructions marked \ours\ are applications developed for this commentary, with explicit literature checks and no priority claim. The audit and commentary are AI-produced by GPT-6 Astra (OpenAI, ChatGPT Pro mode), checked and integrated by Claude (Anthropic).
\endgroup
